% Original assignment starter using CTAN's assignment class (LPPL).
% This document uses only the class interface and contains catalogue-written
% placeholder content, not copied upstream example prose.
\documentclass[11pt,a4paper]{assignment}
\usepackage[T1]{fontenc}
\usepackage[utf8]{inputenc}
\usepackage{lmodern}
\usepackage{microtype}
\usepackage{amsmath,amssymb}
\usepackage{hyperref}
\coursetitle{Systems Performance and Capacity}
\courselabel{SYS 204}
\exercisesheet{Assignment 01: Measuring a Queueing System}{Queueing fundamentals}
\title{Assignment 01: Measuring a Queueing System}
\student{Student Name \\ Student ID}
\semester{Autumn 2026}
\date{August 29, 2026}

\begin{document}

\begin{problemlist}
\pbitem Instructions
\begin{problem}
Show enough working for another student to reproduce your result. State any
assumption that changes the model. Submit this source together with the PDF.
\end{problem}

\pbitem Warm-up: a stable queue
\begin{problem}
A service desk receives requests at rate $\lambda=18$ per hour and serves them
at rate $\mu=24$ per hour. For an $M/M/1$ model, calculate the utilisation
$\rho$, the expected number in the system $L$, and the expected waiting time
$W$.

\begin{answer}
For a stable queue $\lambda<\mu$, so
\[
  \rho=\frac{\lambda}{\mu}=0.75, \qquad
  L=\frac{\rho}{1-\rho}=3, \qquad
  W=\frac{L}{\lambda}=\frac{1}{6}\text{ hour}=10\text{ minutes}.
\]
Explain in your own words what the utilisation value says about spare
capacity during the busiest part of the day.
\end{answer}
\end{problem}

\pbitem Model critique
\begin{problem}
The $M/M/1$ model assumes independent arrivals, exponentially distributed
service times, one server and an unlimited waiting room. Pick two assumptions
that are unlikely to hold for a real software service. For each, describe one
measurement that would test the assumption and one design response if it fails.

\begin{enumerate}
  \item Assumption: \rule{0.62\linewidth}{0.4pt}\\
        Measurement: \rule{0.59\linewidth}{0.4pt}\\
        Response: \rule{0.64\linewidth}{0.4pt}
  \item Assumption: \rule{0.62\linewidth}{0.4pt}\\
        Measurement: \rule{0.59\linewidth}{0.4pt}\\
        Response: \rule{0.64\linewidth}{0.4pt}
\end{enumerate}
\end{problem}

\pbitem Reflection
\begin{problem}
What evidence would convince you that this queueing model is useful enough
for capacity planning, even if it is not literally true?
\end{problem}
\end{problemlist}
\end{document}
